\documentclass[10pt]{beamer}
\usetheme{metropolis}

\usepackage{amsmath,amssymb}
\usepackage{mathtools}
\usepackage{xeCJK}
\setCJKmainfont{Noto Serif CJK SC}
\usepackage{tikz-cd}

% 1. 开启标题背景色填充
\metroset{block=fill}

% theorem environments
\newtheorem{proposition}{Proposition}
\newtheorem{exercise}{Exercise}

% notation
\newcommand{\R}{\mathbb{R}}
\newcommand{\N}{\mathbb{N}}
\newcommand{\F}{\mathcal{F}}
\newcommand{\supp}{\operatorname{supp}}

\title{Ch2\quad Tangent Vectors and Differentials}
\subtitle{Germs of Functions \& The Tangent Space}
\date{September 28, 2026}
\author{LXJ}

\begin{document}

\maketitle

%====================================================================
\section{The Idea: Derivations}

\begin{frame}{Tangent vectors as operators}
	\textbf{Idea.} Any vector $v = (v_1, \dots, v_n) \in \R^n$ can be thought
	of as an \emph{operator} on differentiable functions. If $f$ is
	differentiable at $p$, then $v$ assigns to $f$ the
	\emph{directional derivative}
	\[
		v(f) = \sum_{i=1}^{n} v_i \, \frac{\partial f}{\partial r_i}(p),
		\qquad \text{which is a real value.}
	\]
	It satisfies
	\begin{enumerate}
		\item \textbf{linearity}\quad
		      $v(\lambda f + \mu g) = \lambda\, v(f) + \mu\, v(g)$;
		\item \textbf{Leibniz rule}\quad
		      $v(f \cdot g) = f(p)\cdot v(g) + v(f)\cdot g(p)$.
	\end{enumerate}
\end{frame}

\begin{frame}{Why germs of functions?}
	The operation of taking derivatives only depends on the
	\emph{local} properties of functions. This is captured by the concept
	of \emph{germs of functions}.
	\begin{block}{Remark}
		A vector $v$ is thus ``defined'' as a linear function on
		differentiable functions, determined by the values of $v(f)$.
	\end{block}
	\begin{definition}[Germs of functions]
		Let $m \in M$. Functions $f$ and $g$ defined on open sets
		containing $m$ are said to have the \emph{same germ at $m$} if they
		agree on some neighborhood of $m$. The equivalence classes of germs
		of $C^\infty$ functions are called \emph{germs}, denoted by
		$\widetilde{\F}_m$, i.e.
		\[
			\widetilde{\F}_m = \bigl\{ [f] \bigm|
			f, g \in [f] \ \text{if } \exists \text{ open set } U \ni m,
			\ \text{s.t.\ } f\big|_U = g\big|_U \bigr\}.
		\]
	\end{definition}
\end{frame}

%====================================================================
\section{The Algebra of Germs}

\begin{frame}{Operations on $\widetilde{\F}_m$}
	The addition, scalar multiplication and multiplication of functions
	induce
	\begin{align*}
		\text{(addition)}   & \quad [f] + [g] = [f + g], \\
		\text{(scalar mult.)} & \quad \lambda [f] = [\lambda f], \\
		\text{(multiplication)} & \quad [f] \cdot [g] = [fg].
	\end{align*}
	\begin{exercise}
		Prove that these are well-defined operations.
	\end{exercise}
	\begin{block}{Remark}
		These operations make $\widetilde{\F}_m$ an \emph{algebra} over
		$\R$. A germ $[f]$ has a well-defined value $f(m)$ at $m$.
	\end{block}
	\begin{definition}[The ideals $F_m$ and $F_m^k$]
		\[
			F_m = \bigl\{ [f] \bigm| f \in \widetilde{\F}_m,\ f(m) = 0 \bigr\}.
		\]
	\end{definition}
\end{frame}

\begin{frame}{Powers of the ideal $F_m$}
	$F_m$ is an \emph{ideal} of $\widetilde{\F}_m$. Let $F_m^k$ be the
	$k$-th power of $F_m$: the ideal of $\widetilde{\F}_m$ consisting of
	all finite linear combinations of $k$-fold products of elements of
	$F_m$,
	\[
		F_m^k = \bigl\{ \text{linear comb.\ of } [f_1] \cdots [f_k],
		\ [f_i] \in F_m \bigr\}.
	\]
	These form a descending sequence of ideals
	\[
		\widetilde{\F}_m \supset F_m \supset F_m^2 \supset F_m^3
		\supset \cdots
	\]
\end{frame}

%====================================================================
\section{The Tangent Space}

\begin{frame}{Tangent vectors as derivations}
	\begin{definition}[Tangent vector]
		A \emph{tangent vector} $v$ at the point $m \in M$ is a
		\emph{linear derivation} of the algebra $\widetilde{\F}_m$, i.e.\
		for $[f], [g] \in \widetilde{\F}_m$ and $\lambda, \mu \in \R$:
		\begin{enumerate}
			\item[(a)] $v\bigl(\lambda [f] + \mu [g]\bigr)
				      = \lambda\, v([f]) + \mu\, v([g])$;
			\item[(b)] $v\bigl([f]\cdot[g]\bigr)
				      = v([f])\, g(m) + f(m)\, v([g])$.
		\end{enumerate}
	\end{definition}
	Let $T_m M$ or $M_m$ denote the set of all linear derivations at $m$,
	called the \emph{tangent space to $M$ at $m$}.
	\begin{block}{Note}
		If $M_m \neq \varnothing$, then $M_m$ is a \emph{linear space},
		since for $v, w \in M_m$, $\lambda, \mu \in \R$,
		\[
			(\lambda v + \mu w)([f])
			\ \triangleq\ \lambda\, v([f]) + \mu\, w([f]).
		\]
	\end{block}
\end{frame}

\begin{frame}{Two remarks on the definition}
	\begin{enumerate}
		\item This definition of tangent spaces only works for the
		      $C^\infty$ case;
		\item It is \emph{intrinsic} --- it does not depend on the local
		      coordinate system.
	\end{enumerate}
	\medskip
	\begin{block}{Note (constants are killed)}
		For a derivation $w \in M_m$ and $[c]$ the germ of a locally
		constant function, $w([c]) = 0$.
	\end{block}
	\begin{proof}
		It suffices to show $w([1]) = 0$:
		\[
			w([1]) = w([1]\cdot[1])
			= w([1]) \cdot 1 + 1 \cdot w([1]) = 2\, w([1]). \qedhere
		\]
	\end{proof}
\end{frame}

%====================================================================
\section{$M_m \cong (F_m / F_m^2)^*$}

\begin{frame}{The key isomorphism}
	\begin{lemma}
		$M_m$ is isomorphic to $(F_m / F_m^2)^*$.
	\end{lemma}
	\begin{proof}
		Consider $v \in M_m$. Then $v$ is a linear function on $F_m$ by
		$v : [f] \mapsto v([f])$. This function \emph{vanishes on $F_m^2$}
		by the Leibniz rule of $v$: $\forall\, [f], [g] \in F_m$,
		\[
			v([f][g]) = v([f]) \cdot g(m) + f(m) \cdot v([g]) = 0.
		\]
		So $v$ can be thought of as a linear function defined on the
		cosets $[f] + F_m^2$, therefore $v \in (F_m/F_m^2)^*$.
	\end{proof}
\end{frame}

\begin{frame}{Converse direction}
	Conversely, for $\ell \in (F_m / F_m^2)^*$, we can define
	$v_\ell \in M_m$ by
	\[
		v_\ell([f]) \ \triangleq\
		\ell\bigl( [f] - [f(m)] + F_m^2 \bigr).
	\]
	Then one verifies that $v_\ell$ is a derivation on
	$\widetilde{\F}_m$:
	\begin{enumerate}
		\item \textbf{Linearity.}
		      $v_\ell\bigl(\lambda[f] + \mu[g]\bigr)
		      = \ell\bigl( \lambda[f] - \lambda[f(m)] + F_m^2 \bigr)
		      + \ell\bigl( \mu[g] - \mu[g(m)] + F_m^2 \bigr)
		      = \lambda\, v_\ell([f]) + \mu\, v_\ell([g])$.
		\item \textbf{Leibniz.}
		      $v_\ell\bigl([f][g]\bigr)
		      = \ell\bigl( [f][g] - [f(m)][g(m)] + F_m^2 \bigr)$,
		      where
		      \begin{align*}
			      [f][g] - [f(m)][g(m)]
			      & = \underbrace{\bigl([f]-[f(m)]\bigr)\bigl([g]-[g(m)]\bigr)}_{\in\, F_m^2} \\
			      & \quad + [f(m)]\bigl([g]-[g(m)]\bigr)
			      + \bigl([f]-[f(m)]\bigr)[g(m)],
		      \end{align*}
		      so $v_\ell([f][g]) = f(m)\, v_\ell([g]) + v_\ell([f])\, g(m)$.
	\end{enumerate}
\end{frame}

\begin{frame}{An exercise hidden in the proof}
	\begin{exercise}
		Prove that for $\ell \in (F_m / F_m^2)^*$,
		\[
			\ell\bigl( [f(m)]\cdot\bigl([g]-[g(m)]\bigr) \bigr)
			= f(m)\, \ell\bigl( [g] - [g(m)] \bigr).
		\]
	\end{exercise}
	\bigskip
	\begin{alertblock}{Theorem}
		\[
			\dim\bigl(F_m / F_m^2\bigr) = \dim M.
		\]
	\end{alertblock}
\end{frame}

%====================================================================
\section{Proof of $\dim(F_m/F_m^2) = \dim M$}

\begin{frame}{A Taylor lemma with integral remainder}
	\small
	\begin{lemma}
		If $g \in C^k(U)$ for $k \geq 2$ and $U$ is a convex open set in
		$\R^d$, then $\forall\, p, q \in U$,
		\begin{multline*}
			g(q) = g(p) + \sum_{i=1}^{d}
			\left. \frac{\partial g}{\partial r_i} \right|_{p}
			\bigl(r_i(q) - r_i(p)\bigr) \\
			+ \sum_{i,j} \bigl(r_i(q) - r_i(p)\bigr)
			\bigl(r_j(q) - r_j(p)\bigr)
			\int_0^1 (1-t)
			\left. \frac{\partial^2 g}{\partial r_j \partial r_i}
			\right|_{p + t(q - p)} dt .
		\end{multline*}
	\end{lemma}
	\begin{proof}
		Since $U$ is convex, for any $p, q \in U$ the segment
		$\{ p + t(q-p) \mid t \in [0,1] \} \subset U$. Writing
		$p_i \triangleq r_i(p)$, $q_i \triangleq r_i(q)$, we get
		\[
			g(q) - g(p)
			= \int_0^1 \frac{d}{dt}\, g\bigl(p + t(q-p)\bigr)\, dt
			= \sum_i \int_0^1
			\left. \frac{\partial g}{\partial r_i}
			\right|_{p + t(q-p)} (q_i - p_i)\, dt .
			\qedhere
		\]
	\end{proof}
\end{frame}

\begin{frame}{Proof of the Taylor lemma (continued)}
	For each integrand, expand around $p$:
	\[
		\left. \frac{\partial g}{\partial r_i}
		\right|_{p + t(q-p)}
		= \left. \frac{\partial g}{\partial r_i} \right|_{p}
		+ \sum_j \int_0^t
		\frac{\partial^2 g}{\partial r_j \partial r_i}
		\bigl(p + s(q-p)\bigr)\, (q_j - p_j)\, ds .
	\]
	So, interchanging the order of integration ($s \leftrightarrow t$),
	\begin{align*}
		g(q) & = g(p) + \sum_i
		\left. \frac{\partial g}{\partial r_i} \right|_{p} (q_i - p_i) \\
		& \quad + \sum_{i,j} (q_i - p_i)(q_j - p_j) \int_0^1 (1-t)
		\left. \frac{\partial^2 g}{\partial r_j \partial r_i}
		\right|_{p + t(q-p)} dt .
		\qedhere
	\end{align*}
\end{frame}

\begin{frame}{Proof of the theorem: spanning set}
	Let $(U, \varphi)$ be any coordinate system about $m \in M$, with
	coordinate functions $(x_1, \dots, x_d)$, $d = \dim M$. Let
	$[f] \in \widetilde{\F}_m$. Applying $f \circ \varphi^{-1}$ to the
	lemma:
	\begin{multline*}
		f \circ \varphi^{-1}(x) - f \circ \varphi^{-1}(x_0)
		= \sum_{i=1}^{d}
		\left. \frac{\partial (f \circ \varphi^{-1})}{\partial x_i}
		\right|_{x_0} (x_i - x_i^0) \\
		+ \sum_{i,j} (x_i - x_i^0)(x_j - x_j^0) \int_0^1 (1-t)
		\left. \frac{\partial^2 (f \circ \varphi^{-1})}
		{\partial x_j \partial x_i} \right|_{x_0 + t(x - x_0)} dt .
	\end{multline*}
	Since $x = \varphi(p)$, $x_0 = \varphi(m)$, composing with $\varphi$
	gives, \emph{modulo} $F_m^2$,
	\[
		[f(p)] - [f(m)]
		= \sum_{i=1}^{d}
		\left. \frac{\partial (f \circ \varphi^{-1})}{\partial x_i}
		\right|_{\varphi(m)} \bigl(x_i(p) - x_i(m)\bigr).
	\]
	Hence $\bigl\{ [x_i] - [x_i(m)] \bigr\}$ \emph{spans} $F_m / F_m^2$.
\end{frame}

\begin{frame}{Proof of the theorem: linear independence}
	We claim they are also \emph{linearly independent}. Suppose that
	\[
		\sum_i a_i \bigl(x_i - x_i(m)\bigr) \in F_m^2
		\ \Longleftrightarrow\
		\sum_i a_i \bigl(r_i - r_i(\varphi(m))\bigr) \in F_{\varphi(m)}^2 .
	\]
	Then applying $\left. \dfrac{\partial}{\partial r_j}
	\right|_{\varphi(m)}$ to both sides gives
	\[
		a_j = 0 \qquad \forall\, j .
	\]
	Therefore $\bigl\{ [x_i] - [x_i(m)] \bigr\}_{i=1}^{d}$ is a basis of
	$F_m / F_m^2$, and
	\[
		\dim\bigl(F_m / F_m^2\bigr) = d = \dim M . \qed
	\]
\end{frame}

\end{document}
